\scanpage{130}
\mspage{6}
% archive foliation: 129
pour ton théorème une formulation plus naturelle (qui devrait \deleted{\unsure{en}} fournir une démonstration uniforme, sans distinction du cas $\ell\ne p$ et $\ell=p$), de même qu'on voit en même temps que ta somme abracadabrante se comporte bien quand on fait varier $X$ dans une famille, i.e si $X$ est \deleted{en fait un} un schéma abélien sur $Y$ (car alors les ${}_nX$ seront d'excellents schémas en groupes finis et plats sur $Y$, dont on peut prendre formellement la limite projective\ldots).

J'espère \unsure{arriver-dans|arriver dans} l'année prochaine à une théorie satisfaisante du groupe fondamental, \deleted{\unsure{ce qui}} et achever la rédaction des Chapitres IV,V,VI,VII (ce dernier étant le groupe fondamental), en même temps que des catégories. Dans deux ans résidus, dualité, intersections, Chern, Riemann-Roch. Dans trois ans cohomologie de Weil, et un peu d'homotopie si Dieu veut. Et entre-temps, je ne sais quand, le grand théorème d'existence avec Picard etc, un peu de courbes algébriques, \inserted{les schémas abéliens.} Sauf difficultés imprévues ou enlisement, le multiplodoque devrait être fini d'ici 3 ans, ou 4 au maximum. On pourra commencer à faire de la géométrie algébrique !

Bien à toi
